\documentclass[10pt,a4paper]{article}

% Paquetes para idioma y codificación
\usepackage[utf8]{inputenc}
\usepackage[spanish]{babel}

% Configuración de márgenes súper reducidos para el machete
\usepackage[top=1cm, bottom=1cm, left=1cm, right=1cm]{geometry}

% Paquetes para estructura y formato
\usepackage{multicol} % Para dividir en columnas
\usepackage{enumitem} % Para achicar espacios en las listas
\usepackage{titlesec} % Para achicar los títulos

% Ajuste de tamaño de los títulos para ahorrar espacio
\titleformat{\section}{\large\bfseries}{\thesection}{1em}{}
\titleformat{\subsection}{\normalsize\bfseries}{\thesubsection}{1em}{}
\titlespacing*{\section}{0pt}{2ex}{1ex}
\titlespacing*{\subsection}{0pt}{1ex}{0.5ex}

\begin{document}

% Título del machete
\begin{center}
    \textbf{\Large Machete: Metodologías Ágiles, Scrum y Six Sigma}
\end{center}
\vspace{-0.5cm}
\rule{\textwidth}{0.4pt}
\vspace{0.2cm}

% Inicio de las dos columnas
\begin{multicols}{2}

\section*{1. Metodología SCRUM}
Scrum es uno de los marcos de metodología ágil más populares, que ayuda a los equipos a abordar proyectos complejos dividiendo el trabajo en ciclos iterativos más pequeños llamados sprints. Está diseñado para impulsar la colaboración, aumentar la transparencia e impulsar la mejora continua.

Se suele pensar que scrum y la metodología ágil son lo mismo porque scrum se centra en la mejora continua, que es un principio básico de la metodología ágil. Sin embargo, scrum es un marco para realizar el trabajo, mientras que la metodología ágil es una mentalidad. 

Scrum se basa en varios principios fundamentales que guían su implementación:

\begin{itemize}[leftmargin=*, noitemsep, topsep=0pt]
    \item \textbf{Transparencia:} Todos los aspectos del proceso deben ser visibles para aquellos responsables del resultado. Esto incluye la visibilidad de los trabajos en progreso, así como las reglas y normas que rigen el proyecto.
    \item \textbf{Inspección:} Los elementos del proceso Scrum son revisados regularmente para detectar posibles desviaciones o problemas.
    \item \textbf{Adaptación:} Si durante la inspección se detectan desviaciones significativas, el proceso debe ser ajustado lo antes posible.
    \item \textbf{Iteraciones rápidas:} A través de los sprints, Scrum promueve la entrega rápida de valor, permitiendo al equipo adaptarse a las necesidades cambiantes del proyecto y del cliente. Cada sprint culmina con una revisión y retrospectiva, donde el equipo evalúa tanto el producto desarrollado como el proceso seguido, buscando oportunidades de mejora continua.
\end{itemize}

Dentro de la metodología Scrum, los roles definidos son esenciales para el éxito del proyecto. Cada uno tiene responsabilidades específicas que aseguran que el equipo funcione de manera eficiente y que el producto final cumpla con las expectativas del cliente. Estos roles no solo facilitan la organización interna, sino que también promueven la colaboración y la transparencia:

\begin{itemize}[leftmargin=*, noitemsep, topsep=0pt]
    \item \textbf{Scrum Master:} Su papel es asegurar que el equipo siga los principios y prácticas de Scrum correctamente, eliminando cualquier obstáculo que pueda impedir el progreso del equipo. Este rol no es el de un jefe tradicional, sino el de un líder servidor que apoya al equipo en alcanzar su máximo potencial.
    \item \textbf{Product Owner:} Es el responsable de maximizar el valor del producto que está siendo desarrollado. Este rol actúa como el puente entre los stakeholders (interesados) y el equipo de desarrollo, asegurando que las necesidades del cliente se traduzcan en características del producto.
    \item \textbf{Equipo de Desarrollo:} En Scrum es un grupo multifuncional y autoorganizado que es responsable de entregar un incremento de producto al final de cada sprint. Este equipo está compuesto por profesionales con todas las habilidades necesarias para completar el trabajo de desarrollo.
    \item \textbf{Stakeholders:} Son las partes interesadas en el proyecto, que pueden incluir clientes, usuarios finales, patrocinadores, y otros actores externos que tienen un interés en el resultado del producto.
\end{itemize}

\section*{2. White Belt}
El White Belt Six Sigma es una certificación básica dentro de la metodología Six Sigma que introduce a los participantes en conceptos clave como la reducción de variabilidad, mejora de procesos y enfoque en el cliente. A diferencia de los niveles Green o Black Belt, los White Belts no lideran proyectos, sino que colaboran apoyando actividades y entendiendo el lenguaje común de Six Sigma en la organización.

Incorporar formación White Belt en tu organización puede traer múltiples ventajas que fortalecen tanto los procesos como la cultura empresarial:

\begin{itemize}[leftmargin=*, noitemsep, topsep=0pt]
    \item \textbf{Conciencia generalizada:} Al capacitar a más empleados con conocimientos básicos, se genera un lenguaje común y un entendimiento colectivo sobre la mejora de calidad.
    \item \textbf{Incremento en la colaboración:} Los White Belts pueden apoyar proyectos liderados por cinturones superiores, lo que mejora la sinergia y el trabajo en equipo.
    \item \textbf{Detección temprana de problemas:} Con más personas capacitadas, es más fácil identificar oportunidades de mejora o desviaciones en procesos.
    \item \textbf{Reducción de costos:} La mejora continua desde la base ayuda a minimizar desperdicios, errores y retrabajos que afectan el presupuesto.
    \item \textbf{Compromiso y motivación:} Involucrar a colaboradores en iniciativas de calidad genera mayor sentido de pertenencia y responsabilidad.
\end{itemize}

Para sacar el máximo provecho de esta certificación básica, considera estos pasos:

\begin{enumerate}[leftmargin=*, noitemsep, topsep=0pt]
    \item Identificar a los colaboradores clave: Selecciona empleados de diferentes áreas para que reciban la capacitación y sirvan de enlace con sus equipos.
    \item Seleccionar un programa de formación reconocido: Opta por cursos que cubran los conceptos fundamentales y que sean prácticos y accesibles.
    \item Fomentar la aplicación práctica: Motiva a los White Belts a participar en pequeños proyectos o actividades de mejora continua.
    \item Integrar a niveles superiores: Asegurar que White Belts trabajen coordinadamente con Green y Black Belts para impulsar resultados efectivos.
    \item Medir resultados: Llevar un seguimiento de los beneficios obtenidos, como reducción de defectos o mejoras en tiempos de proceso.
\end{enumerate}

El White Belt Six Sigma es una herramienta valiosa para involucrar a toda la organización en la mejora continua, permitiendo a tu empresa ser más competitiva, eficiente y enfocada en la calidad.

\section*{3. Green Belt}
Un Green Belt es alguien certificado en Lean Six Sigma con conocimiento intermedio. No es un experto estadístico, pero sí domina las herramientas necesarias para identificar problemas, liderar proyectos de mejora y tomar decisiones basadas en datos.

Los Green Belts son empleados de una organización que han sido capacitados en la metodología de mejora Six Sigma y dirigirán un equipo de mejora de procesos como parte de su trabajo a tiempo completo. Dedican más tiempo a los componentes de toma de decisiones y creación de estrategias de los elementos clave en el proceso de planificación del proyecto Six Sigma. Tienen una comprensión más profunda del proceso general; trabajan con gerentes de proyectos centralizados en la entrega de retroalimentación y la conducción de los objetivos de rendimiento.

Así pues, los profesionales que obtienen el certificado de Six Sigma Green Belt, se preparan para ser expertos en las siguientes funciones:

\begin{itemize}[leftmargin=*, noitemsep, topsep=0pt]
    \item Liderar proyectos Kaizen o de mejora de procesos.
    \item Aplicar la metodología DMAIC.
    \item Facilitar sesiones de resolución de problemas.
    \item Colaborar con Black Belts en proyectos más grandes.
    \item Analizar datos de rendimiento y proponer soluciones efectivas.
\end{itemize}

\section*{4. Metodología Ágil}
El esquema o metodología agile es una metodología iterativa, es decir, se realizan entregas cíclicas y en cada entrega se realizan todas las fases del ciclo: desde toma de requerimientos, diseño, verificación y entrega. La mayor diferencia de las metodologías ágiles, frente a los antiguos modelos waterfall, es que en los procesos ágiles se entrega valor y recibe feedback constantemente durante todo el proyecto. En lugar de depender de un producto final predefinido desde el inicio, estas metodologías promueven la creación de software funcional que evoluciona de manera incremental.

Las fases del proyecto, como el desarrollo, diseño y prueba, se realizan en ciclos cortos, conocidos como sprints, permitiendo que el equipo se ajuste rápidamente a los cambios y reciba feedback de los clientes de forma continua. Después de cada sprint, los equipos reflexionan y observan lo que ha sucedido. Evalúan si hay algo que se podría mejorar para poder ajustar la estrategia para el siguiente sprint.

Los principios básicos son:

\begin{itemize}[leftmargin=*, noitemsep, topsep=0pt]
    \item \textbf{Individuos e interacciones sobre procesos y herramientas.} Valorar más a los individuos y sus interacciones que a los procesos y las herramientas. Los equipos autoorganizados que trabajan con metodologías ágiles valoran más la colaboración en equipo y trabajar juntos, que trabajar de manera independiente y hacer las cosas “al pie de la letra”.
    \item \textbf{Software funcionando sobre documentación extensiva.} Valorar más el software en funcionamiento que la documentación exhaustiva. El software que desarrollan los equipos ágiles debe funcionar.
    \item \textbf{Colaboración con el cliente sobre negociación contractual.} Valorar más la colaboración entre las partes interesadas y el cliente que la negociación contractual.
    \item \textbf{Respuesta ante el cambio sobre seguir un plan.} Valorar más la respuesta ante el cambio que seguir un plan. Uno de los principales beneficios de las metodologías ágiles de proyectos es que permiten que los equipos sean flexibles.
\end{itemize}

Los 4 valores son los pilares de la metodología agile. A partir de esos valores, el equipo desarrolló 12 principios. Estos principios se pueden adaptar fácilmente para adecuarlos a las necesidades de cada equipo. Si los valores ágiles fueran los pilares fundamentales de una casa, los 12 principios serían las habitaciones que se pueden construir dentro:

\begin{enumerate}[leftmargin=*, noitemsep, topsep=0pt]
    \item Satisfacer a los clientes a través de la entrega temprana y continua.
    \item Aceptamos que los requisitos cambien, incluso en etapas tardías del proyecto.
    \item Hacemos entregas valiosas con frecuencia.
    \item Rompemos con el aislamiento en los proyectos. La colaboración es la clave de la estrategia ágil.
    \item Desarrollamos proyectos con personas motivadas.
    \item El modo más eficiente de comunicar es con conversaciones cara a cara.
    \item El software en funcionamiento es el principal indicador del progreso. Lo más importante es que los equipos se esfuercen (con la aplicación de estos esquemas de trabajo ágiles) por desarrollar los productos.
    \item Mantenemos un ritmo de trabajo constante.
    \item La excelencia continua favorece la agilidad.
    \item La simplicidad es esencial. A veces, la solución más simple es la mejor.
    \item La organización interna y autónoma de los equipos genera resultados más valiosos.
    \item Con regularidad, el equipo reflexiona y adapta las formas de trabajo para favorecer la efectividad.
\end{enumerate}

\end{multicols}

\end{document}